% ============================================================
\subsection{Recursion and Memoization}
% ============================================================

A recursive function is defined by three obligations.

\begin{keybox}[The recursion contract]
\begin{enumerate}[itemsep=1ex]
  \item \textbf{Base case.} At least one input for which the answer is
        returned without a recursive call.
  \item \textbf{Recursive case.} The answer for input $n$ expressed in terms
        of the answers for strictly smaller inputs.
  \item \textbf{Progress.} Every recursive call moves the argument closer to a
        base case, so the recursion terminates.
\end{enumerate}

\medskip

Correctness follows by induction. Assuming the recursive call returns the
correct answer on the smaller input, verify only that the combination step is
correct. Tracing the full call tree by hand is unnecessary.
\end{keybox}

\paragraph{The call stack.} Each pending call occupies a stack frame, so the
space cost of a recursion is $O(d)$ where $d$ is the maximum depth. Python
imposes a default recursion limit of 1000 frames and raises
\texttt{RecursionError} beyond it. A recursion whose depth scales with $n$ is
therefore unsafe for large $n$ unless it is rewritten as a loop with an
explicit stack.

\paragraph{Overlapping subproblems.} Naive recursion recomputes identical
subproblems. The Fibonacci recurrence $F(n) = F(n-1) + F(n-2)$ generates a
call tree with $O(\varphi^n)$ nodes, since $F(n-2)$ is evaluated once by the
$F(n-1)$ branch and again directly. Caching each result by its argument
collapses the tree to $O(n)$ distinct evaluations.

\begin{lstlisting}[style=python, caption={Memoization by decorator}]
from functools import lru_cache

@cache(maxsize=None)     
def fib(n):
    if n < 2:                     # base case
        return n
    return fib(n-1) + fib(n-2)    # recursive case
\end{lstlisting}

Arguments to a memoized function must be hashable, which is the reason
memoized state is passed as an integer, a string, or a tuple rather than as a
list or a set.

\paragraph{Top-down and bottom-up.} Memoized recursion is the top-down form.
It computes only the states that are reachable, at the cost of a call stack.
Tabulation is the bottom-up form. It fills an array in an order that
guarantees every dependency is already resolved, using no stack and no
recursion.

\begin{lstlisting}[style=python, caption={The same recurrence, tabulated}]
def fib(n):
    if n < 2:
        return n
    prev, curr = 0, 1
    for _ in range(n - 1):
        prev, curr = curr, prev + curr  # only two states are ever live
    return curr
\end{lstlisting}

\vspace{-2ex}

The final line illustrates a general reduction. When a recurrence depends only
on the previous $k$ states, the table collapses from $O(n)$ space to $O(k)$.

% ------------------------------------------------------------
\subsubsection*{Dynamic programming}

Dynamic programming applies when a problem has \textbf{optimal substructure},
meaning the optimal solution is composed of optimal solutions to subproblems,
together with \textbf{overlapping subproblems}, meaning those subproblems recur.
Absent overlap, plain divide and conquer suffices. Absent optimal substructure,
a greedy or exhaustive method is required.

\begin{keybox}[Formulating a DP]
Five questions, answered in order:
\vspace{0.5ex}

\begin{enumerate}[itemsep=1ex]
  \item \textbf{State.} What arguments fully determine a subproblem? This is
        the signature of the function or the index set of the table.
  \item \textbf{Recurrence.} How is a state's value expressed in terms of
        smaller states?
  \item \textbf{Base cases.} Which states are known outright?
  \item \textbf{Order.} In what order can states be evaluated so that
        dependencies are resolved first? Top-down recursion answers this
        automatically.
  \item \textbf{Answer.} Which state holds the final result?
\end{enumerate}

\medskip

The complexity is the number of states multiplied by the cost of evaluating
one state.
\end{keybox}

\paragraph{Canonical state definitions.} Most problems reuse a small number of
state shapes.

\begin{center}
\begin{tabular}{lll}
\toprule
\textbf{State} & \textbf{Meaning} & \textbf{Typical problem} \\
\midrule
\texttt{dp[i]}    & best over the prefix ending at $i$ & maximum subarray, house robber \\
\texttt{dp[i][j]} & best over two prefixes             & edit distance, longest common subsequence \\
\texttt{dp[i][w]} & first $i$ items, capacity $w$      & 0/1 knapsack \\
\texttt{dp[i][j]} & best over the interval $[i,j]$     & matrix chain, palindrome partition \\
\bottomrule
\end{tabular}
\end{center}

% ------------------------------------------------------------
\subsubsection*{Backtracking}

Backtracking is recursion over a search tree with an explicit undo step. It
enumerates permutations, subsets, and combinations, and prunes branches that
cannot yield a valid solution.

\newpage

\begin{lstlisting}[style=python, caption={Backtracking skeleton}]
def solve(path, choices, results):
    if is_complete(path):
        results.append(path[:]) # copy, since path is mutated in place
        return
    for c in choices:
        if not is_valid(path, c):
            continue                   # prune
        path.append(c)                 # choose
        solve(path, remaining(choices, c), results)
        path.pop()                     # undo
\end{lstlisting}

The \texttt{path[:]} copy is mandatory. Appending \texttt{path} itself stores a
reference to a list that subsequent mutation will alter, so every recorded
result ends up identical.

For unconstrained enumeration, the standard library is preferable to a
hand-written recursion.

\begin{lstlisting}[style=python, caption={Enumeration without recursion}]
from itertools import permutations, combinations, product

permutations(xs, r)     # ordered arrangements of length r
combinations(xs, r)     # unordered selections of length r
product(xs, repeat=r)   # r-tuples with repetition, the r-fold Cartesian power
\end{lstlisting}

These correspond exactly to the counting formulas $P(n,r) = n!/(n-r)!$,
$\binom{n}{r}$, and $n^r$.
